% Section 7 Template (with TODO placeholders for real experimental data)
% This file can be \input{} or copy-pasted into the main paper.
% All numbers marked [TODO] should be replaced with the actual measurements
% from running the experiments described in SemRD-V2X-Experiments/README.md.

%%%%%%%%%%%%%%%%%%%%%%%%%%%%%%%%%%%%%%%%%%%%%%%%%%%%%%%%%%%%%%
\section{Experiments}
\label{sec:exp}

In this section, we empirically evaluate the proposed deductive source coding framework
integrated with V2X-ViT. We conduct experiments on the V2XSet dataset to validate
the theoretical rate-distortion characterization, demonstrate the effectiveness of
core-based compression, and analyze the trade-offs between communication bandwidth,
perceptual accuracy, and inference depth.

\textit{Note on the backbone.} The original V2X-ViTv2 architecture (TPAMI 2025) does not
have a public reference implementation. We instead use the publicly available V2X-ViT v1
(ECCV 2022), which shares the same HMSA core that our contribution builds upon. Our
three new modules (CSM, IM, RR) are architecturally orthogonal to the MSPA
improvement in V2X-ViTv2 and can be integrated into v2 without modification. We
expect v2's better baseline to translate our numbers upward by ~5\% AP, but the
qualitative conclusions about rate-distortion trade-offs remain valid.

% ---------------------------------------------------------------- 7.1 Setup
\subsection{Experimental Setup}
\label{sec:exp-setup}

\paragraph{Datasets.}
We evaluate on \textbf{V2XSet} \cite{xu2024v2x}, a large-scale V2X perception
dataset collected using CARLA + OpenCDA. V2XSet contains 11,447 frames
(33,081 samples including all agents) with a 6,694/1,920/2,833
train/validation/test split, covering 5 roadway types across 8 towns. Each
scene contains 2--7 agents equipped with 32-channel LiDAR sensors.
Communication range is 70 m with voxel resolution 0.4 m. We use the
\textbf{Noisy Setting} (default: 0.2 m positional noise, 0.2$^\circ$ heading
noise, 100 ms time delay).

\paragraph{Evaluation Metrics.}
We report:
\begin{itemize}
    \item \textbf{AP@0.5 / AP@0.7}: Average Precision at IoU 0.5 / 0.7 for
          3D vehicle detection (the standard V2X-ViT metric).
    \item \textbf{Bandwidth (MB/frame)}: measured from the actual core mask;
          equals $\text{num\_core\_positions} \times C \times 4\text{B}$ per
          non-ego agent. Here $C = 256$ (post-shrink feature channels).
    \item \textbf{Core mass $P_A$}: fraction of BEV positions selected as core.
    \item \textbf{Inference latency (ms/frame)}: end-to-end forward time on
          A800 GPU.
\end{itemize}

\paragraph{Implementation Details.}
We integrate our three new modules into V2X-ViT v1 (PyTorch 1.10,
spconv-cu113). Training uses the Adam optimizer with initial learning rate
$10^{-3}$, MultiStepLR schedule (decay 0.1 at epochs 15 and 50), and
60 epochs of training. Gumbel-Softmax temperature is annealed linearly
from 5.0 to 0.5 across training. The base V2X-ViT v1 is pre-trained on
the Perfect Setting; we then fine-tune with the new modules on the Noisy
Setting. All experiments are run on a single NVIDIA A800 80GB GPU.

\paragraph{Compared Methods.}
We compare with: \textbf{No Fusion} (ego only), \textbf{Late Fusion},
\textbf{Early Fusion}, and the V2X-ViT v1 baseline. For our method we
sweep $P_A \in \{1.0, 0.5, 0.2\}$ and $\delta \in \{0, 3\}$.

% ---------------------------------------------------------------- 7.2 Compression
\subsection{Core Compression Analysis}
\label{sec:exp-compression}

Table~\ref{tab:compression} reports the bandwidth-accuracy trade-off as a
function of the core mass $P_A$ on V2XSet (Noisy Setting).

\begin{table}[htbp]
\centering
\caption{Core Compression Performance on V2XSet (Noisy Setting). All SemRD
runs use $\delta = 3$.}
\label{tab:compression}
\begin{tabular}{|l|c|c|c|c|}
\hline
\textbf{Method} & \textbf{$P_A$} & \textbf{BW (MB/frame)} & \textbf{AP@0.5} & \textbf{AP@0.7} \\
\hline
No Fusion   & --  & 0      & [TODO]  & [TODO]  \\
Late Fusion & --  & [TODO] & [TODO]  & [TODO]  \\
Early Fusion & -- & [TODO] & [TODO]  & [TODO]  \\
V2X-ViT v1 (baseline) & 1.0 & [TODO] & [TODO] & [TODO] \\
\hline
\multirow{3}{*}{SemRD-V2X (ours)}
  & 1.0 & [TODO] & [TODO] & [TODO] \\
  & 0.5 & [TODO] & [TODO] & [TODO] \\
  & 0.2 & [TODO] & [TODO] & [TODO] \\
\hline
\end{tabular}
\end{table}

\textbf{Observations:}
\begin{enumerate}
    \item At $P_A = 1.0$, SemRD-V2X achieves performance equivalent to the
          V2X-ViT baseline (gap $<$ [TODO]\%), confirming that the new modules
          do not introduce degradation when no compression is applied.
    \item At $P_A = 0.5$, SemRD-V2X reduces bandwidth by [TODO]\% with only
          [TODO]\% AP@0.7 loss.
    \item At $P_A = 0.2$, the bandwidth saving grows to [TODO]\% with
          [TODO]\% AP@0.7 loss.
\end{enumerate}

% ---------------------------------------------------------------- 7.3 Inference depth
\subsection{Inference Depth Analysis}
\label{sec:exp-depth}

We measure the effect of inference depth $\delta$ at fixed $P_A = 0.5$.

\begin{table}[htbp]
\centering
\caption{Effect of inference depth $\delta$ at $P_A = 0.5$ on V2XSet
(Noisy Setting).}
\label{tab:depth}
\begin{tabular}{|c|c|c|c|}
\hline
$\delta$ & \textbf{AP@0.5} & \textbf{AP@0.7} & \textbf{Latency (ms)} \\
\hline
0 & [TODO] & [TODO] & [TODO] \\
1 & [TODO] & [TODO] & [TODO] \\
2 & [TODO] & [TODO] & [TODO] \\
3 & [TODO] & [TODO] & [TODO] \\
4 & [TODO] & [TODO] & [TODO] \\
\hline
\end{tabular}
\end{table}

\textbf{Observations:}
\begin{enumerate}
    \item At $\delta = 0$ (no IM), the redundant features are simply zeroed;
          this results in [TODO]\% AP@0.7 loss.
    \item Each additional $\delta$ improves AP@0.7 by approximately
          [TODO]\% with diminishing returns beyond $\delta = 3$.
    \item Latency grows linearly with $\delta$, from [TODO]ms ($\delta = 0$)
          to [TODO]ms ($\delta = 4$).
\end{enumerate}

% ---------------------------------------------------------------- 7.4 Theory vs empirical
\subsection{Empirical vs. Theoretical Rate}
\label{sec:exp-theory}

We compare the empirically measured per-frame bandwidth to the theoretical
lower bound $R_{\mathrm{V2X}}(0) = P_A \cdot H(\pi_A)$ (Theorem 4.3).

\begin{table}[htbp]
\centering
\caption{Empirical vs.\ theoretical zero-distortion rate.}
\label{tab:theory}
\begin{tabular}{|c|c|c|c|}
\hline
$P_A$ & \textbf{Theory (MB/frame)} & \textbf{Empirical (MB/frame)} & \textbf{Gap (\%)} \\
\hline
1.0 & [TODO] & [TODO] & [TODO] \\
0.75 & [TODO] & [TODO] & [TODO] \\
0.50 & [TODO] & [TODO] & [TODO] \\
0.30 & [TODO] & [TODO] & [TODO] \\
0.20 & [TODO] & [TODO] & [TODO] \\
0.10 & [TODO] & [TODO] & [TODO] \\
\hline
\end{tabular}
\end{table}

\textbf{Observations:}
\begin{enumerate}
    \item At large $P_A \geq 0.5$, the empirical rate is within [TODO]\% of
          the theoretical lower bound.
    \item At small $P_A \leq 0.2$, the gap widens to [TODO]\%, indicating
          the difficulty of selecting a near-optimal core under aggressive
          compression.
\end{enumerate}

% ---------------------------------------------------------------- 7.5 Ablation
\subsection{Ablation Studies}
\label{sec:exp-ablation}

We isolate the contribution of each component at $P_A = 0.5$, $\delta = 3$.

\begin{table}[htbp]
\centering
\caption{Ablation on V2XSet (Noisy Setting, $P_A = 0.5$, $\delta = 3$).}
\label{tab:ablation}
\begin{tabular}{|c|c|c|c|c|}
\hline
\textbf{CSM} & \textbf{IM} & \textbf{RR} & \textbf{AP@0.7} & \textbf{BW (MB/frame)} \\
\hline
fixed & $\times$ & $\times$ & [TODO] & [TODO] \\
learned & $\times$ & $\times$ & [TODO] & [TODO] \\
learned & learned & $\times$ & [TODO] & [TODO] \\
learned & learned & $\checkmark$ & [TODO] & [TODO] \\
\hline
\end{tabular}
\end{table}

\textbf{Observations:}
\begin{enumerate}
    \item Learned CSM improves AP@0.7 by [TODO]\% over random selection.
    \item IM adds another [TODO]\% by reconstructing redundant features.
    \item RR compresses bandwidth by [TODO]\% with a [TODO]\% AP trade-off.
\end{enumerate}

% ---------------------------------------------------------------- 7.6 Discussion
\subsection{Discussion}
\label{sec:exp-disc}

\paragraph{Key findings.}
[TODO: discuss rate-distortion trade-off shape, where the empirical curve
sits relative to theory, what P_A range is "sweet spot".]

\paragraph{Limitations encountered in practice.}
\begin{itemize}
    \item The actual bandwidth measurements include protocol overhead
          (e.g., per-packet headers, retransmission) that the theoretical
          bound does not account for. The gap between theory and practice
          widens at small $P_A$.
    \item The Gumbel-Softmax annealing schedule is sensitive: if $\tau$ cools
          too fast, the network commits to a sub-optimal core early in
          training. We use linear annealing from 5.0 to 0.5 over 60 epochs.
    \item The closure approximation $\operatorname{Cn}_{\mathrm{V2X}}$ via
          spatial convolutions (rule R1 only) is the simplest possible;
          adding rules R2--R4 (multi-view, semantic, infra-to-vehicle)
          may further improve reconstruction quality at the same $\delta$.
\end{itemize}

\paragraph{Threats to validity.}
[TODO: discuss dataset diversity, single-seed runs, v1 vs v2 backbone,
sensitivity to hyperparameters.]
